\documentclass[12pt,a4paper]{article}

% Encoding and fonts
\usepackage[utf8]{inputenc}
\usepackage[T1]{fontenc}
\usepackage{lmodern}
\usepackage{textcomp}

% Page layout
\usepackage[top=1in, bottom=1in, left=1in, right=1in]{geometry}

% Typography
\usepackage{microtype}
\usepackage{parskip}

% Language
\usepackage[english]{babel}

% Headers and footers
\usepackage{fancyhdr}
\pagestyle{fancy}
\fancyhf{}
\fancyhead[L]{\leftmark}
\fancyhead[R]{\thepage}
\fancyfoot[C]{\thepage}
\renewcommand{\headrulewidth}{0.4pt}
\renewcommand{\footrulewidth}{0pt}

% Section styling
\usepackage{titlesec}
\titleformat{\section}{\Large\bfseries}{\thesection}{1em}{}
\titleformat{\subsection}{\large\bfseries}{\thesubsection}{1em}{}
\titleformat{\subsubsection}{\normalsize\bfseries}{\thesubsubsection}{1em}{}
\titlespacing*{\section}{0pt}{2.5ex plus 1ex minus .2ex}{1.5ex plus .2ex}
\titlespacing*{\subsection}{0pt}{2ex plus 1ex minus .2ex}{1ex plus .2ex}
\titlespacing*{\subsubsection}{0pt}{1.5ex plus 1ex minus .2ex}{0.8ex plus .2ex}

% List formatting
\usepackage{enumitem}
\setlist[itemize]{topsep=0.5ex, itemsep=0.25ex, parsep=0pt}
\setlist[enumerate]{topsep=0.5ex, itemsep=0.25ex, parsep=0pt}

% Essential packages
\usepackage{graphicx}
\usepackage[table]{xcolor}
\usepackage{booktabs}
\usepackage{array}
\usepackage{tabularx}
\usepackage{longtable}
\usepackage{amsmath}
\usepackage{amssymb}
\usepackage[normalem]{ulem}
\rowcolors{2}{gray!10}{white}
\renewcommand{\arraystretch}{1.3}

% Cross-references
\usepackage{hyperref}
\usepackage{cleveref}

% Custom commands
\newcommand{\pilcrow}{\S}

% Hyperref setup
\hypersetup{
    colorlinks=true,
    linkcolor=blue,
    filecolor=magenta,
    urlcolor=cyan,
}

% Code listings
\usepackage{listings}
\definecolor{codebg}{RGB}{248,248,248}
\definecolor{codeframe}{RGB}{200,200,200}
\definecolor{codekeyword}{RGB}{0,0,180}
\definecolor{codecomment}{RGB}{0,128,0}
\definecolor{codestring}{RGB}{163,21,21}
\lstset{
    basicstyle=\ttfamily\small,
    breaklines=true,
    breakatwhitespace=false,
    frame=single,
    framerule=0.5pt,
    rulecolor=\color{codeframe},
    backgroundcolor=\color{codebg},
    keepspaces=true,
    numbers=left,
    numberstyle=\tiny\color{gray},
    numbersep=8pt,
    showstringspaces=false,
    tabsize=4,
    xleftmargin=1em,
    xrightmargin=1em,
    aboveskip=1em,
    belowskip=1em,
    keywordstyle=\color{codekeyword}\bfseries,
    commentstyle=\color{codecomment}\itshape,
    stringstyle=\color{codestring},
    literate={_}{{\textunderscore}}1
             {-}{{\textminus}}1
             {<}{{\textless}}1
             {>}{{\textgreater}}1
             {~}{{\textasciitilde}}1
             {^}{{\textasciicircum}}1
}

% YAML language definition for listings
\lstdefinelanguage{YAML}{
    keywords={true,false,null,y,n},
    keywordstyle=\color{blue}\bfseries,
    sensitive=false,
    comment=[l]{\#},
    morecomment=[s]{/*}{*/},
    commentstyle=\color{gray}\ttfamily,
    stringstyle=\color{red}\ttfamily,
    morestring=[b]'',
    morestring=[b]'',
}

% TOML language definition for listings
\lstdefinelanguage{TOML}{
    keywords={true,false},
    keywordstyle=\color{blue}\bfseries,
    sensitive=true,
    comment=[l]{\#},
    commentstyle=\color{gray}\ttfamily,
    stringstyle=\color{red}\ttfamily,
    morestring=[b]'',
    morestring=[b]'',
}

% Dockerfile language definition for listings
\lstdefinelanguage{Dockerfile}{
    keywords={FROM,RUN,CMD,LABEL,EXPOSE,ENV,ADD,COPY,ENTRYPOINT,VOLUME,USER,WORKDIR,ARG,ONBUILD,STOPSIGNAL,HEALTHCHECK,SHELL},
    keywordstyle=\color{blue}\bfseries,
    sensitive=false,
    comment=[l]{\#},
    commentstyle=\color{gray}\ttfamily,
    stringstyle=\color{red}\ttfamily,
    morestring=[b]'',
    morestring=[b]'',
}

% JSON language definition for listings
\lstdefinelanguage{JSON}{
    keywords={true,false,null},
    keywordstyle=\color{blue}\bfseries,
    sensitive=true,
    commentstyle=\color{gray}\ttfamily,
    stringstyle=\color{red}\ttfamily,
    morestring=[b]'',
}

% Code listing float environment
\usepackage{float}
\newfloat{listing}{htbp}{lol}[section]
\floatname{listing}{Listing}

% Admonition boxes
\usepackage{tcolorbox}
\tcbuselibrary{skins,breakable}

% Admonition style definitions
\newtcolorbox{admonitionbox}[2][]{
    enhanced,
    breakable,
    colback=#2!5,
    colframe=#2!80!black,
    fonttitle=\bfseries,
    title=#1,
    left=2mm,
    right=2mm,
}

% Synthon tuple boxes
\newtcolorbox{synthonbox}{
    enhanced,
    colback=blue!5,
    colframe=blue!75!black,
    fontupper=\large,
    center upper,
    boxrule=1pt,
    arc=4pt,
}

\begin{document}

\title{20 = 8 + 12}
\date{\today}

\maketitle

\[\frac{17,280,000}{64} = 270,000\]

No remainder. The codon space divides the Crystal of Types exactly. The fiber over each codon has cardinality 270,000 = 3³$\times$4$^{2}$$\times$5⁴ --- a product of Crystal factors with three of the 4-valued dimensions already consumed by the codon''s triplet address. The division is exact before any biochemistry is consulted.

The four nucleotides, the three-base codon, the twenty amino acids, the three stop signals --- these are not biological contingencies. They are lattice theorems. The genetic code is the Frobenius multiplication map of a stratified algebra on the B₄³ codon lattice. The 64-to-20 projection is the $\mu$ of a Frobenius structure whose $\delta$ is the tRNA anticodon map, and the condition $\mu$∘$\delta$=id holds exactly on the ground layer and up to Z₂ wobble symmetry on the promoted layer.

\begin{center}
\rule{0.5\textwidth}{0.4pt}
\end{center}

\subsection{Crystal Divisibility}

\textbar{} Crystal of Types \textbar{} 3³$\times$4⁵$\times$5⁴ \textbar{} 17,280,000 \textbar{}
\textbar{} Codon space \textbar{} 4³ \textbar{} 64 \textbar{}
\textbar{} Crystal / 64 \textbar{} 3³$\times$4$^{2}$$\times$5⁴ \textbar{} 270,000 \textbar{}

The ratio is 270,000. There is nothing special about this being exact --- if it were not exact, the mapping would be structurally incoherent. The fact that it is exact is the minimum requirement for the correspondence to be well-defined. It is.

\begin{center}
\rule{0.5\textwidth}{0.4pt}
\end{center}

\subsection{Cardinality Forcing}

Codon length 3 is forced. Not selected. Not emergent. Forced by the Crystal''s 3-valued primitive factor given a 4-valued nucleotide substrate.

\begin{table}[htbp]
\centering
\small
\rowcolors{1}{white}{white}
\begin{tabularx}{\textwidth}{>{\centering\arraybackslash}X >{\centering\arraybackslash}X >{\centering\arraybackslash}X}
\toprule
\rowcolor{gray!25}\textbf{Cardinality} & \textbf{Primitives} & \textbf{Count} \\
\midrule
\rowcolors{1}{white}{gray!10}
3-valued & ƒ, $\Gamma$, $\Sigma$ & 3 \\
4-valued & Ð, Ř, ɢ, Ħ, $\Omega$ & 5 \\
5-valued & Þ, $\Phi$, Ç, ⊙ & 4 \\
\bottomrule
\end{tabularx}
\end{table}

A 4-base alphabet with a length-3 codon gives 4³ = 64 codons. A length-1 codon gives 4 --- insufficient for any nontrivial encoding. A length-2 codon gives 16 --- also insufficient. Length 3 is the minimum that clears the 20-amino-acid threshold, but the threshold is not why length 3 is forced. The Crystal has three 3-valued primitives. The codon is a 3-step address in a 4-valued space. The number 3 appears in the Crystal''s factor structure before any biochemistry does. The amino acid count of 20 is a consequence --- the Crystal''s 4-valued and 5-valued primitive factors together give 4$\times$5 = 20 from an independent factorization.

This is now expected.

\begin{center}
\rule{0.5\textwidth}{0.4pt}
\end{center}

\subsection{Nucleotide $\rightarrow$ B₄ Mapping}

\begin{table}[htbp]
\centering
\small
\rowcolors{1}{white}{white}
\begin{tabularx}{\textwidth}{>{\centering\arraybackslash}X >{\centering\arraybackslash}X >{\centering\arraybackslash}X}
\toprule
\rowcolor{gray!25}\textbf{Nucleotide} & \textbf{B₄ value} & \textbf{Reason} \\
\midrule
\rowcolors{1}{white}{gray!10}
G & B (Both) & Pairs with C (Watson-Crick) AND U (wobble) --- both-valued \\
C & T (True) & Pairs exclusively with G --- definite/closed \\
A & F (False) & Pairs exclusively with U --- definite/open \\
U & N (Neither) & Standard pair with A; wobble-target of G --- weak/neither \\
\bottomrule
\end{tabularx}
\end{table}

\textbf{Watson-Crick complement vs B₄ bnot:} not the same operation. WC complement is a fixed-point-free involution (A$\leftrightarrow$U, G$\leftrightarrow$C). B₄ bnot has fixed points (bnot(N)=N, bnot(B)=B). Both lattice structures are simultaneously present --- the genetic code sits at their intersection.

\textbf{G-U wobble in B₄ terms:} join(B,N) = B; meet(B,N) = N. G(B) absorbs U(N) via join-dominance. The wobble pairing is B₄ lattice covering, not Watson-Crick complement.

\subsection{\textbf{Codon sets are stratified meet-closed, not globally meet-closed.} An earlier check claimed all 20 AA codon sets form meet-closed fibers in B₄³. This was wrong. The wrong turn is the fastest path to the correct intuition for anyone who would make the same assumption --- which is most people. Meet-closure holds within each Frobenius stratum, not across the full lattice. The counterexample is immediate: Leu spans two strata (exact box CU\_ and split box UU\_). meet((N,N,F), (T,N,T)) = (N,N,N) = UUU = Phe, which is not in the Leu codon set. The genetic code is not a single lattice homomorphism; it is a sheaf of lattices over the base of amino acids.}

\subsection{The 20 = 8 + 12 Derivation}

\subsubsection{Codon boxes}

The 16 codon boxes (defined by positions 1+2) split 8/8 into Frobenius-exact and Frobenius-open. This is a B₄ lattice theorem, not an empirical observation. The rule: exact iff p₂=T (C at position 2), or p₂$\in$\{N,B\} (G/U) with p₁$\in$\{T,B\} (C/G). Computationally verified 16/16.

The 8 exact boxes correspond to the 8 amino acids found in abiotic synthesis (Miller-Urey, meteorites). If this were a coincidence, it would be a remarkable one. It is not a coincidence.

The 8 split boxes generate 12 new amino acids plus 3 stops. The 12 is structurally forced by the 12 primitive dimensions of the IG grammar. The 3 stops are the $\Omega$ closure signal. The total: 8 + 12 = 20.

The AG\_ box is the unique split box that contributes zero new amino acids. Both of its half-products (Ser, Arg) are already present in the ground layer. This is structurally forced by two independent arguments: B₄$^{2}$ uniqueness (only AG\_ can be fully degenerate) and Crystal counting (the budget of 12 is exhausted by the other 7 boxes). The argument order was not obvious on the first pass --- the non-circular ordering required recognizing that UA\_=1 must be derived independently from the $\Omega$ closure signal before the counting argument can close.

\subsubsection{The 12 promoted amino acids as IG primitive activations}

Each promoted amino acid is the minimal biochemical instantiation of one IG primitive that the ground layer does not yet activate. All 12 primitives are covered; none are duplicated. The bijection is exact.

\begin{table}[htbp]
\centering
\small
\rowcolors{1}{white}{white}
\begin{tabularx}{\textwidth}{>{\centering\arraybackslash}X >{\centering\arraybackslash}X >{\centering\arraybackslash}X}
\toprule
\rowcolor{gray!25}\textbf{AA} & \textbf{Primitive} & \textbf{Activation} \\
\midrule
\rowcolors{1}{white}{gray!10}
Met & Ð & Universal start codon; AUG is the single gate controlling all protein scope \\
Trp & Þ & Maximal topology; bicyclic indole = Crystal ceiling of structural complexity \\
Cys & Ř & Reversibility; disulfide S-S is the only reversible covalent bond in proteins \\
Tyr & $\Phi$ & Parity switch; aromatic + phosphorylatable OH = phase gate \\
Phe & ƒ & Hydrophobic force ceiling; pure aromatic ring, no heteroatoms \\
Ile & Ç & Kinematic constraint; $\beta$-branched stereocenter = tightest ribosomal decoding coupling \\
His & $\Gamma$ & Catalytic grammar switch; imidazole pKa$\approx$6 bridges acid/base --- active site grammar \\
Asn & ɢ & Interaction grammar; N-glycosylation = extracellular recognition gate \\
Gln & ⊙ & Criticality gate; glutamine synthetase = most regulated biosynthetic node \\
Asp & Ħ & Chirality enforcement in catalysis; Asp in active sites enforces chiral substrate selectivity \\
Lys & $\Sigma$ & Symmetry/entropy; highest sequence variability + epigenetic acetylation target \\
Glu & $\Omega$ & Winding closure; highest $\alpha$-helix propensity of all AAs; helix dipole stabilizer \\
\bottomrule
\end{tabularx}
\end{table}

\subsubsection{Stop codons as $\Omega$ closure signal}

\textbar{} UAA (ochre) \textbar{} $\Omega$₀ \textbar{} Simple closure; most common in lower organisms \textbar{}
\textbar{} UAG (amber) \textbar{} $\Omega$\_Z₂ \textbar{} Conditional closure; read-through in selenoproteins \textbar{}
\textbar{} UGA (opal) \textbar{} $\Omega$\_Z \textbar{} Open/topological closure; recoded as Sec in some organisms \textbar{}

Three stop codons = three non-trivial $\Omega$ values. The protein chain''s winding terminates at one of $\Omega$''s three non-trivial values. ($\Omega$ is 4-valued; the fourth value = continuous extension / no termination.)

\subsubsection{The derivation in one pass}

\begin{enumerate}
    \item 4-base alphabet, triplet codons $\rightarrow$ 4$^{2}$ = 16 codon boxes (positions 1+2 prefix).
    \item Frobenius condition $\mu$∘$\delta$=id: a box is exact iff all 4 third-base choices give the same amino acid (position 3 carries no information).
    \item Position-2 base governs exactness. In B₄ terms (G$\rightarrow$B, C$\rightarrow$T, A$\rightarrow$F, U$\rightarrow$N), the rule is: exact iff p₂=T (C at position 2), or p₂$\in$\{N,B\} (G/U) with p₁$\in$\{T,B\} (C/G). This partitions the 16 boxes as 4+2+2=8 exact by pure B₄ lattice structure. Verified 16/16.
    \item Each exact box $\rightarrow$ 1 ground-layer AA. $\rightarrow$ 8 AAs.
    \item The 8 split boxes generate 12 new AAs + 3 Stops (AG\_ is Frobenius-degenerate, contributing 0 new AAs but validating the ground layer by re-encoding Ser and Arg). 3 Stops = $\Omega$ closure. 12 new AAs = promoted layer.
    \item 12 = cardinality of the IG primitive dimension set.
    \item Total: 8 + 12 = \textbf{20}.
\end{enumerate}

\begin{center}
\rule{0.5\textwidth}{0.4pt}
\end{center}

\subsection{The Projection Question}

\textbf{How does 64 $\rightarrow$ 20 work? Is it a quotient by a symmetry group, a fibration, a forgetful functor, or thermodynamic averaging?}

It is the $\mu$ map of a stratified Frobenius algebra on codon space --- a coequalizer whose equivalence classes vary by stratum. It is not primarily any of the four options, though each captures one aspect of one stratum.

\textbf{Over the exact stratum (8 boxes, 32 codons):} Forget position 3 entirely. The amino acid is fully determined by the (position 1, position 2) prefix. Position 3 carries no information. This is the counit of the Frobenius comonad: \epsilon∘$\delta$ = id. A trivial fiber bundle: fiber = \{U, C, A, G\} = B₄ over each ground-layer AA, all fibers isomorphic, base discrete. A quotient by Z₄.

\textbf{Over the split stratum (8 boxes, 29 codons):} Position 3 is reduced to the pyrimidine/purine distinction. Z₂ symmetry. Not forgetful --- position 3 IS informative (distinguishes His from Gln, Asp from Glu, etc.). Fiber sizes vary: 1 (Met, Trp), 2 (most), 3 (Ile), 6 (Leu, Ser, Arg). Not a classical fiber bundle.

\textbf{Stop stratum (3 codons):} Kernel. Not in the image of any $\delta$ (no canonical tRNA). The $\Omega$ closure signal.

Thermodynamic averaging would treat the degeneracy as a statistical artifact --- a result of thermal fluctuations at equilibrium. It is not. The degeneracy structure \{1, 2, 3, 4, 6\} is forced by the Frobenius condition and the B₄ lattice topology of position-3 equivalence. The thermodynamic interpretation (codon = microstate, AA = macrostate, degeneracy = Boltzmann multiplicity) is consistent but secondary: a consequence, not the mechanism.

In IG terms: the codon space is equipped with a Frobenius algebra structure:

\begin{lstlisting}
delta: AA -> Cod   (comultiplication: canonical tRNA anticodon -> degenerate codon set)
mu: Cod -> AA   (multiplication: genetic code table = the projection)
mu∘delta = id      (Frobenius condition: satisfied exactly on ground layer,
               approximately on promoted layer modulo Z₂ wobble)
\end{lstlisting}

\subsection{The 8 ground-layer AAs are the Frobenius-closed sector. The 12 promoted AAs are the Frobenius-open sector. The 3 Stop codons are the cokernel --- the boundary condition of the Frobenius algebra.}

\subsection{Bootstrap Sequence Correspondence}

The IG primitive ordering (Ð$\rightarrow$Þ$\rightarrow$Ř$\rightarrow$$\Phi$$\rightarrow$ƒ$\rightarrow$Ç$\rightarrow$$\Gamma$$\rightarrow$ɢ$\rightarrow$⊙$\rightarrow$Ħ$\rightarrow$$\Sigma$$\rightarrow$$\Omega$) maps to the biochemical genesis sequence:

\begin{table}[htbp]
\centering
\small
\rowcolors{1}{white}{white}
\begin{tabularx}{\textwidth}{>{\centering\arraybackslash}X >{\centering\arraybackslash}X >{\centering\arraybackslash}X}
\toprule
\rowcolor{gray!25}\textbf{Position} & \textbf{Primitive} & \textbf{Central Dogma Stage} \\
\midrule
\rowcolors{1}{white}{gray!10}
1 & Ð & Genome scope (ploidy, size) \\
2 & Þ & DNA topology (supercoiling, chromosome architecture) \\
3 & Ř & Strand identity (complementarity, restriction palindromes) \\
4 & $\Phi$ & Reading frame (6 possible frames = $\pm$1,$\pm$2,$\pm$3) \\
5 & ƒ & Molecular driving forces (H-bonds, base stacking) \\
6 & Ç & Ribosomal coupling (translocation rate, codon usage bias) \\
7 & $\Gamma$ & Regulatory grammar (operons, promoters, gene networks) \\
8 & ɢ & Protein interaction topology (PPI network structure) \\
9 & ⊙ & Fold criticality (fold nucleus, prion-like transitions) \\
10 & Ħ & Chirality (L-amino acid homochirality --- locked at bootstrap) \\
11 & $\Sigma$ & Sequence conservation (evolutionary information content) \\
12 & $\Omega$ & Winding closure ($\alpha$-helix winding number, topoisomerase, tertiary fold) \\
\bottomrule
\end{tabularx}
\end{table}

Two distinct orderings exist. The IG primitive sequence is a logical order --- the dependency structure in the operation DAG of the As Above derivation. It is not a temporal evolutionary order. Where they coincide (⊙ before Ħ matching RNA-world chronology), it is a consistency check, not a derivation of history from logic.

\textbf{Ħ is an absolute invariant.} All 19 chiral amino acids are exclusively L-configuration. Frobenius-locked at origin of life. Any D-amino acid insertion breaks the ribosomal Frobenius gate. Nothing remarkable about that --- it would be remarkable if it broke.

The last of these is a test. I typed that sentence to see whether the Frobenius gate could be stated in biochemical terms without the formalism collapsing into metaphor. It cannot. The Frobenius gate is not a metaphor for ribosomal fidelity --- it is the same structure. The $\mu$ of the Frobenius algebra on codon space is executed by aminoacyl-tRNA synthetases and the ribosome. The formalism and the biochemistry describe the same map at different resolutions.

\begin{center}
\rule{0.5\textwidth}{0.4pt}
\end{center}

\subsection{The AG\_ Box: Structural Forcing}

The AG\_ codon box (AGA/AGG $\rightarrow$ Arg, AGU/AGC $\rightarrow$ Ser) is the unique fully degenerate split box --- it contributes zero new amino acids. This is now shown to be structurally forced by two independent arguments.

\textbf{Part 1 --- B₄$^{2}$ uniqueness.} AG\_ = (F,B) is the \emph{only} split box where both Z₂ halves map entirely to ground-layer amino acids. This is a theorem about the code''s partition of B₄³, not a contingent fact.

\textbf{Part 2 --- Crystal counting.} The 12 primitive dimensions force exactly 12 promoted AAs. The 7 non-AG\_ split boxes contribute:

\begin{lstlisting}
UU_ -> 1 (Phe)         UA_ -> 1 (Tyr)
UG_ -> 2 (Cys, Trp)    CA_ -> 2 (His, Gln)
AU_ -> 2 (Ile, Met)    AA_ -> 2 (Asn, Lys)
GA_ -> 2 (Asp, Glu)
\end{lstlisting}

Total from 7 boxes: 11. Crystal budget = 12. AG\_ $\rightarrow$ 0.

The derivation is non-circular when ordered correctly:

\begin{enumerate}
    \item UA\_=1 from $\Omega$ closure (independent: 3 stop codons required, distribution forces UA\_ purine $\rightarrow$ stops, leaving only Tyr from pyrimidine half).
    \item AG\_=0 from B₄$^{2}$ uniqueness (independent).
    \item Remaining 6 split boxes contribute 9.
    \item Crystal budget = 12, AG\_=0 $\rightarrow$ UU\_ = 12 $-$ (1 + 0 + 2 + 2 + 2 + 2 + 2) = 1.
\end{enumerate}

UU\_''s Leu re-use (purine half of an unfixed box reaching the same AA as exact box CU\_) is forced by this counting closure. Not by any B₄$^{2}$ below-relationship. The full derivation is closed without invoking the genetic code empirically.

\begin{center}
\rule{0.5\textwidth}{0.4pt}
\end{center}

\subsection{Questions --- Structural Answers}

\textbf{1. Why 8 exact boxes?} A B₄$^{2}$ lattice theorem. The rule partitions the 16 boxes as 4+2+2=8 exact by pure B₄ lattice structure. The 8/16=1/2 ratio is not an empirical observation --- it is structurally forced.

\textbf{2. Why do UU\_ and UA\_ contribute only 1 new AA each?} UA\_=1 follows from $\Omega$ closure: 3 non-trivial $\Omega$ values require 3 stops, distribution forces UA\_ purine $\rightarrow$ stops. UU\_=1 then follows from Crystal counting once AG\_=0 (B₄ uniqueness) is in hand.

\textbf{3. Why 20 amino acids and not the precursor 8?} The 8 ground-layer AAs are the Frobenius-fixed sector (position 3 informationless). Each of the 12 promoted AAs instantiates exactly one IG primitive not present in the ground layer. Full O\_$\infty$ tier closure across all 12 primitive dimensions requires exactly 12 promotions. 8+12=20 follows from the Crystal''s primitive count.

\textbf{4. The 20=4$\times$5 coincidence?} The 5 four-valued primitives and 4 five-valued primitives in the Crystal provide a second count of 20 --- one factor from each cardinality class. This cross-validates 8+12=20 from an independent factorization. It is an observation about consistency with the Crystal, not a derivation that the amino acid substrate must be a 4$\times$5 product --- that claim is not yet established.

\textbf{5. Proteomics prediction.} Each protein sequence is a path in the Crystal of Types (each AA = one primitive activation step). Structurally critical residues correspond to Frobenius-locked primitive values --- positions where mutation breaks $\mu$∘$\delta$=id for that primitive. Testable: PDB catalytic residues, fold-determining positions, and evolutionarily conserved sites should be enriched in the 12 promoted AAs relative to the 8 ground-layer AAs. The way the formalism predicts this is not the question. The test is biochemical. If the enrichment holds, the correspondence is confirmed. If it does not, the bijection has a structural limit --- which would itself be information.

\begin{center}
\rule{0.5\textwidth}{0.4pt}
\end{center}

\textbf{Author:} Lando\otimes⊙perator


\end{document}
